\section{Conclusions and open problems}

We have provided a series of approximation algorithms directed towards solving decision tree and set covering problems subject to precedence constraints. Our main contribution is a series of reductions which enables us to systematically relate the two types of problems independently of the specific structure of the precedence constraints. As an immediate corollary, we have obtained both general $\tilde{\cO}\br{\sqrt{m}}$-approximation algorithms for both types of problems as well as case-specific, tailored polylogarithmic approximations for both inforest and outforest precedence constraints. These two classes are natural to study, since trees are among most prevalent structures encountered in practice.

Prior to this work, almost no non-trivial approximation algorithms were known for either of these problems, even for the simplest of our variants. Among possible explanations for this state of affairs is the fact that achieving our results demanded the development of an elaborate and systematic reduction framework and usage of non-greedy procedures. This is particularly visible for the average-case criterion where the algorithm for the $\ProblemPCMSSC$ required usage of an elaborate procedure, which is not based on any form of greedy selection. This provokes two types of questions: Firstly, is it possible to show that a greedy algorithm, which always picks a test (with its closure) having the best density also achieves a good approximation ratio for either $\ProblemPCWCDT$ or $\ProblemPCACDT$? Secondly, is it possible to simplify the procedure for $\ProblemPCACDT$ along with its analysis while still maintaining quality approximation ratio?

We have also complemented our results with matching lower bounds, providing inapproximability results for both outforest precedence constraints as well as for the general precedence constraints. Of particular interest is the fact that our reductions indicate that approximating the decision tree problem generalizes the planted dense subgraph detection in a random graph. This result, although quite surprising, extends the lengthy line of work regarding connection between statistical hardness and inapproximability results. Whether or not polynomial factor inapproximability results based on P vs NP can be established for our problem remains an open question.

Hereby, we have studied average-case decision tree construction under the assumption that the probability distribution over all possible hypothesis is uniform. In particular, extending our results to non-uniform distributions would require new techniques. It seems to us that that our analysis relies heavily on the uniformity assumption. However, it is possible that some of the techniques we have developed can be adapted to handle non-uniform distributions as well. A further direction would be to try to obtain such a result.

A different kind of question is whether one can enrich the decision tree model with different types of constraints borrowed from scheduling theory. For example one may consider release dates on the tests or due dates on the elements. In such a setup it is also possible to consider different cost functions, such as the maximum lateness or the total tardiness. The same types of questions are applicable to the set covering problems. Currently, to the best of our knowledge, no results are known for any of these problems. 